\section{Descripción del problema}
\label{cap:problema}

Queremos responder preguntas sobre los hábitos de escucha de los usuarios de ListenBrainz: cuántas escuchas se procesan por día y cómo varían, cuánto escucha en promedio un usuario activo, qué tan concentrada está la escucha en pocos artistas, qué artistas tienen más alcance y qué artistas comparten público. Además, el flujo tiene que poder recibir datos nuevos de forma periódica, porque ListenBrainz publica incrementales y la consigna pide simular actualizaciones.

Desde el punto de vista de los datos, estas preguntas no tienen la misma forma:

\begin{itemize}
  \item Conservar el registro original de cada escucha, con campos opcionales y metadatos que cambian de un registro a otro, es un problema de documentos semiestructurados.
  \item Contar escuchas por día, calcular crecimiento y armar rankings mensuales sobre millones de filas es un problema de series temporales con consultas conocidas de antemano.
  \item Encontrar artistas que comparten oyentes, detectar comunidades o recomendar artistas a un usuario a partir de lo que escuchan usuarios parecidos es un problema de relaciones.
\end{itemize}

Una sola base de datos puede guardar todo, pero obliga a resolver mal alguna de estas partes. Un modelo relacional o documental puede contar escuchas por día, pero cada consulta recorre millones de registros si no se precalcula el agregado. Recorrer relaciones de varios saltos (usuario, artista, otro artista, otros usuarios) en una base documental o columnar exige uniones costosas o reconstruir el grafo en memoria. Y una base de grafos no es la herramienta adecuada para guardar millones de documentos JSON originales ni para servir series temporales.

Por eso planteamos una arquitectura multimodelo en la que cada tipo de dato vive en la base que mejor responde sus consultas, y un orquestador que mantiene las tres sincronizadas a partir de la misma fuente. Hay que tener presente una restricción del dataset durante todo el informe: un incremental de ListenBrainz refleja las escuchas \emph{enviadas} a la plataforma en ese periodo, no toda su historia ni la popularidad general de los artistas. Los resultados describen ese lote de ingesta, no a la población de oyentes.
